\begin{ex}
Lớp 3A có 15 bạn thích màu đỏ, 12 bạn thích màu xanh, 8 bạn thích màu vàng. Hỏi có bao nhiêu bạn thích màu đỏ?
\shortans{15}
\loigiai{
Có 15 bạn thích màu đỏ.
}
\end{ex}

\begin{ex}
Theo số liệu trên, tổng số học sinh thích cả 3 màu là bao nhiêu?
\shortans{35}
\loigiai{
$15 + 12 + 8 = 35$ bạn.
}
\end{ex}

\begin{ex}
Số bạn thích màu đỏ nhiều hơn số bạn thích màu vàng bao nhiêu bạn?
\shortans{7}
\loigiai{
$15 - 8 = 7$ bạn.
}
\end{ex}

\begin{ex}
Trong một cuộc khảo sát về thú cưng: 18 bạn thích chó, 14 bạn thích mèo, 6 bạn thích thỏ. Thú cưng nào được nhiều bạn thích nhất? (Chó ghi 1, mèo ghi 2, thỏ ghi 3).
\shortans{1}
\loigiai{
Chó có 18 bạn thích (nhiều nhất).
}
\end{ex}

\begin{ex}
Tổng số bạn tham gia khảo sát thú cưng trên là bao nhiêu?
\shortans{38}
\loigiai{
$18 + 14 + 6 = 38$ bạn.
}
\end{ex}

\begin{ex}
Một tuần bạn Lan ghi chép số trang sách đọc được mỗi ngày: Thứ Hai: 10, Thứ Ba: 12, Thứ Tư: 15, Thứ Năm: 8, Thứ Sáu: 15. Ngày nào Lan đọc được ít trang sách nhất? (Ghi số thứ, ví dụ Thứ Năm ghi 5).
\shortans{5}
\loigiai{
Thứ Năm (8 trang, ít nhất).
}
\end{ex}

\begin{ex}
Tổng số trang sách Lan đọc được từ Thứ Hai đến Thứ Sáu là bao nhiêu trang?
\shortans{60}
\loigiai{
$10 + 12 + 15 + 8 + 15 = 60$ trang.
}
\end{ex}

\begin{ex}
Số trang sách Lan đọc nhiều nhất trong một ngày là bao nhiêu?
\shortans{15}
\loigiai{
Nhiều nhất là 15 trang.
}
\end{ex}

\begin{ex}
Bảng theo dõi số cây trồng của 4 tổ: Tổ 1: 25 cây, Tổ 2: 30 cây, Tổ 3: 28 cây, Tổ 4: 32 cây. Tổ trồng nhiều cây nhất là tổ mấy?
\shortans{4}
\loigiai{
Tổ 4 trồng được 32 cây.
}
\end{ex}

\begin{ex}
Tổng số cây cả 4 tổ trồng được là bao nhiêu cây?
\shortans{115}
\loigiai{
$25 + 30 + 28 + 32 = 115$ cây.
}
\end{ex}

\begin{ex}
Tổ 4 trồng nhiều hơn Tổ 1 bao nhiêu cây?
\shortans{7}
\loigiai{
$32 - 25 = 7$ cây.
}
\end{ex}

\begin{ex}
Bảng ghi số học sinh giỏi của 3 lớp: Lớp 3A: 12 bạn, Lớp 3B: 15 bạn, Lớp 3C: 10 bạn. Cả 3 lớp có bao nhiêu học sinh giỏi?
\shortans{37}
\loigiai{
$12 + 15 + 10 = 37$ bạn.
}
\end{ex}

\begin{ex}
Lớp có số học sinh giỏi ít nhất có bao nhiêu bạn?
\shortans{10}
\loigiai{
Lớp 3C có 10 bạn.
}
\end{ex}

\begin{ex}
Số học sinh giỏi lớp 3B nhiều hơn lớp 3C bao nhiêu bạn?
\shortans{5}
\loigiai{
$15 - 10 = 5$ bạn.
}
\end{ex}

\begin{ex}
Số lượng bóng các màu trong rổ: Đỏ: 8 quả, Xanh: 12 quả, Vàng: 10 quả. Tổng số bóng trong rổ là bao nhiêu quả?
\shortans{30}
\loigiai{
$8 + 12 + 10 = 30$ quả.
}
\end{ex}

\begin{ex}
Số bóng xanh gấp số bóng đỏ bao nhiêu lần? (Viết dưới dạng phân số tối giản hoặc tỉ số: lấy số bóng xanh chia cho 4 được mấy?).
\shortans{3}
\loigiai{
$12 : 4 = 3$.
}
\end{ex}

\begin{ex}
Hiệu giữa số bóng xanh và số bóng đỏ là bao nhiêu quả?
\shortans{4}
\loigiai{
$12 - 8 = 4$ quả.
}
\end{ex}

\begin{ex}
Số bông hoa 4 bạn hái được: Mai: 9 bông, Lan: 12 bông, Cúc: 8 bông, Đào: 11 bông. Ai hái được ít hoa nhất? (Mai ghi 1, Lan ghi 2, Cúc ghi 3, Đào ghi 4).
\shortans{3}
\loigiai{
Cúc hái được 8 bông (ít nhất).
}
\end{ex}

\begin{ex}
Tổng số hoa 4 bạn hái được là bao nhiêu bông?
\shortans{40}
\loigiai{
$9 + 12 + 8 + 11 = 40$ bông.
}
\end{ex}

\begin{ex}
Số điểm 10 của bạn Minh trong tuần: Thứ Hai: 2, Thứ Ba: 3, Thứ Tư: 1, Thứ Năm: 2, Thứ Sáu: 4. Tổng số điểm 10 Minh đạt được là bao nhiêu?
\shortans{12}
\loigiai{
$2 + 3 + 1 + 2 + 4 = 12$ điểm 10.
}
\end{ex}

\begin{ex}
Ngày Minh đạt nhiều điểm 10 nhất có bao nhiêu điểm 10?
\shortans{4}
\loigiai{
Có 4 điểm 10 vào thứ Sáu.
}
\end{ex}

\begin{ex}
Theo bảng số liệu: Khối 1 có 120 học sinh, Khối 2 có 130 học sinh, Khối 3 có 140 học sinh. Cả 3 khối có bao nhiêu học sinh?
\shortans{390}
\loigiai{
$120 + 130 + 140 = 390$ học sinh.
}
\end{ex}

\begin{ex}
Khối 3 nhiều hơn Khối 1 bao nhiêu học sinh?
\shortans{20}
\loigiai{
$140 - 120 = 20$ học sinh.
}
\end{ex}

\begin{ex}
Bảng số liệu bán hàng ngày Chủ nhật: Bán được 45 cái bút, 60 quyển vở, 20 cục tẩy. Tổng số món đồ đã bán là bao nhiêu?
\shortans{125}
\loigiai{
$45 + 60 + 20 = 125$ món.
}
\end{ex}

\begin{ex}
Số quyển vở bán được nhiều hơn số cái bút bao nhiêu cái?
\shortans{15}
\loigiai{
$60 - 45 = 15$.
}
\end{ex}
